\SourceLeaf{348}%
bei konkav-konvexen Flächen das negative. Besteht das betrachtete Stück der gekrümmten Fläche in dieser Hinsicht aus Teilen verschiedener Art, so sind genauere Bestimmungen erforderlich, die hier übergangen werden müssen.

Vergleicht man die Flächeninhalte zweier einander entsprechender Stücke der gekrümmten Fläche und der Oberfläche der Hilfskugel, so gelangt man zu einem neuen Begriff, ähnlich wie aus dem Vergleich von Volumen und Masse der Begriff der Dichte hervorgeht. Als \textit{Krümmungsmaß} in einem Punkt der gekrümmten Fläche bezeichnet der Verfasser den Wert des Bruchs, dessen Nenner der Flächeninhalt eines unendlich kleinen Stücks der gekrümmten Fläche an diesem Punkt ist. Im Zähler steht der Flächeninhalt des entsprechenden Stücks der Hilfskugel, also die ganze Krümmung dieses Elements. Im Sinne des Verfassers entsprechen die ganze Krümmung und das Krümmungsmaß bei gekrümmten Flächen also dem, was man bei gekrümmten Linien als Amplitude beziehungsweise als Krümmung schlechthin bezeichnet. Er hatte Bedenken, die beiden letzteren Ausdrücke auf gekrümmte Flächen zu übertragen, weil sie eher aus Gewohnheit als wegen ihrer Eignung gebräuchlich sind. Auf die Bezeichnungen selbst kommt es weniger an als darauf, dass prägnante Sätze ihre Einführung rechtfertigen.

Die Aufgabe, das Krümmungsmaß in jedem Punkt einer gekrümmten Fläche zu finden, nimmt je nach der Art, wie die Fläche gegeben ist, verschiedene Formen an. Am einfachsten ist es, die Punkte im Raum allgemein durch drei rechtwinklige Koordinaten \(x,y.\ z\) zu kennzeichnen und eine Koordinate als Funktion der beiden anderen darzustellen. So erhält man den einfachsten Ausdruck für das Krümmungsmaß. Zugleich ergibt sich ein bemerkenswerter Zusammenhang zwischen diesem Maß und den Krümmungen der Kurven, die beim Schnitt der gekrümmten Fläche mit Ebenen senkrecht zu ihr entstehen. Bekanntlich hat \textsc{Euler} zuerst gezeigt, dass zwei dieser Schnittebenen, die einander unter einem rechten Winkel schneiden, die Eigenschaft haben, dass in der einen der größte, in der anderen der kleinste Krümmungsradius auftritt oder, genauer gesagt, dass in ihnen die beiden extremen Krümmungen vorkommen. Aus dem erwähnten Ausdruck für das Krümmungsmaß folgt nun, dass es einem Bruch mit dem Zähler eins und dem Produkt der beiden extremen Krümmungsradien als Nenner gleich ist. Weniger einfach ist der Ausdruck für das Krümmungsmaß, wenn
